\documentclass[10pt,a4paper]{amsart}
\usepackage{amsmath,amssymb,amsthm,graphicx}
% Retained preamble for the Stacks Project chapter morphisms.tex (Morphisms of Schemes).
% The chapter file pinned-original-source.tex carries only \input{preamble} on its first line, so no
% preamble text lives inside the chapter itself.  The Stacks Project uses one single shared preamble for
% every chapter; the snapshot stages it verbatim in sources/stacks-chow/preamble.tex.  The block below is
% copied mechanically, line for line, from that file's lines 14-18 (the Xy-pic setup) and lines 154-217
% (the native theorem declarations together with the mathematical macros).  Nothing was added, removed or
% redefined.  The wrapper class and the amsmath/amssymb/amsthm/graphicx base packages are supplied by the
% assembler, so the class/\externaldocument part of the shared preamble is not repeated here.
% For commutative diagrams we use Xy-pic
\usepackage[all]{xy}

% We use 2cell for 2-commutative diagrams.
\xyoption{2cell}
% Theorem environments.
%
\theoremstyle{plain}
\newtheorem{theorem}[subsection]{Theorem}
\newtheorem{proposition}[subsection]{Proposition}
\newtheorem{lemma}[subsection]{Lemma}

\theoremstyle{definition}
\newtheorem{definition}[subsection]{Definition}
\newtheorem{example}[subsection]{Example}
\newtheorem{exercise}[subsection]{Exercise}
\newtheorem{situation}[subsection]{Situation}

\theoremstyle{remark}
\newtheorem{remark}[subsection]{Remark}
\newtheorem{remarks}[subsection]{Remarks}

\numberwithin{equation}{subsection}

% Macros
%
\def\lim{\mathop{\mathrm{lim}}\nolimits}
\def\colim{\mathop{\mathrm{colim}}\nolimits}
\def\Spec{\mathop{\mathrm{Spec}}}
\def\Hom{\mathop{\mathrm{Hom}}\nolimits}
\def\Ext{\mathop{\mathrm{Ext}}\nolimits}
\def\SheafHom{\mathop{\mathcal{H}\!\mathit{om}}\nolimits}
\def\SheafExt{\mathop{\mathcal{E}\!\mathit{xt}}\nolimits}
\def\Sch{\mathit{Sch}}
\def\Mor{\mathop{\mathrm{Mor}}\nolimits}
\def\Ob{\mathop{\mathrm{Ob}}\nolimits}
\def\Sh{\mathop{\mathit{Sh}}\nolimits}
\def\NL{\mathop{N\!L}\nolimits}
\def\CH{\mathop{\mathrm{CH}}\nolimits}
\def\proetale{{pro\text{-}\acute{e}tale}}
\def\etale{{\acute{e}tale}}
\def\QCoh{\mathit{QCoh}}
\def\Ker{\mathop{\mathrm{Ker}}}
\def\Im{\mathop{\mathrm{Im}}}
\def\Coker{\mathop{\mathrm{Coker}}}
\def\Coim{\mathop{\mathrm{Coim}}}

% Boxtimes
%
\DeclareMathSymbol{\boxtimes}{\mathbin}{AMSa}{"02}

%
% Macros for moduli stacks/spaces
%
\def\QCohstack{\mathcal{QC}\!\mathit{oh}}
\def\Cohstack{\mathcal{C}\!\mathit{oh}}
\def\Spacesstack{\mathcal{S}\!\mathit{paces}}
\def\Quotfunctor{\mathrm{Quot}}
\def\Hilbfunctor{\mathrm{Hilb}}
\def\Curvesstack{\mathcal{C}\!\mathit{urves}}
\def\Polarizedstack{\mathcal{P}\!\mathit{olarized}}
\def\Complexesstack{\mathcal{C}\!\mathit{omplexes}}
% \Pic is the operator that assigns to X its picard group, usage \Pic(X)
% \Picardstack_{X/B} denotes the Picard stack of X over B
% \Picardfunctor_{X/B} denotes the Picard functor of X over B
\def\Pic{\mathop{\mathrm{Pic}}\nolimits}
\def\Picardstack{\mathcal{P}\!\mathit{ic}}
\def\Picardfunctor{\mathrm{Pic}}
\def\Deformationcategory{\mathcal{D}\!\mathit{ef}}

% --- appended 2026-10-12 (Chief): the project own definitions of the reference/slogan/history
% display environments, copied mechanically from the shared Stacks preamble lines 8-12
% (\usepackage{verbatim} plus the three \newenvironment lines, which the chapter preamble keeps
% verbatim from the project preamble). Nothing added, removed or redefined: these environments
% carry no mathematical content and the project itself typesets them as comments. Without these
% three lines a counted statement whose own text uses them fails with Environment slogan undefined.
\usepackage{verbatim}
\newenvironment{reference}{\comment}{\endcomment}
\newenvironment{slogan}{\comment}{\endcomment}
\newenvironment{history}{\comment}{\endcomment}
\usepackage{cleveref}
\crefname{originallemma}{original lemma}{original lemmas}
\Crefname{originallemma}{Original lemma}{Original lemmas}
\makeatletter
\newlabel{algebra-lemma-finite-over-subring}{{0560}{source}{}{sourcepointer-1}{}}
\newlabel{algebra-lemma-finite-over-subring@cref}{{[originallemma][0][]0560}{[1][0][]source}{}{}{}}
\makeatother
\crefname{originallemma}{original lemma}{original lemmas}
\Crefname{originallemma}{Original lemma}{Original lemmas}
\makeatletter
\newlabel{algebra-lemma-support-closed}{{00L2}{source}{}{sourcepointer-2}{}}
\newlabel{algebra-lemma-support-closed@cref}{{[originallemma][0][]00L2}{[1][0][]source}{}{}{}}
\makeatother
\crefname{originallemma}{original lemma}{original lemmas}
\Crefname{originallemma}{Original lemma}{Original lemmas}
\makeatletter
\newlabel{lemma-closed-immersion-bijection-ideals}{{01QQ}{source}{}{sourcepointer-3}{}}
\newlabel{lemma-closed-immersion-bijection-ideals@cref}{{[originallemma][0][]01QQ}{[1][0][]source}{}{}{}}
\makeatother
\crefname{originallemma}{original lemma}{original lemmas}
\Crefname{originallemma}{Original lemma}{Original lemmas}
\makeatletter
\newlabel{lemma-i-star-equivalence}{{01QY}{source}{}{sourcepointer-4}{}}
\newlabel{lemma-i-star-equivalence@cref}{{[originallemma][0][]01QY}{[1][0][]source}{}{}{}}
\makeatother
\crefname{originallemma}{original lemma}{original lemmas}
\Crefname{originallemma}{Original lemma}{Original lemmas}
\makeatletter
\newlabel{properties-lemma-finite-type-module}{{01PB}{source}{}{sourcepointer-5}{}}
\newlabel{properties-lemma-finite-type-module@cref}{{[originallemma][0][]01PB}{[1][0][]source}{}{}{}}
\makeatother
\title{Existence, uniqueness and finiteness of the scheme theoretic support of a finite type quasi-coherent module}
\author{The Stacks Project Authors}
\date{stacks-project:morphisms.tex:local-snapshot-20261012; GNU FDL 1.2 or later}
\begin{document}
\maketitle
\noindent\textit{Editorial note.} All delivered mathematical text is the author's own, copied from the pinned Stacks Project chapter file pinned-original-source.tex; no statement or proof step was rewritten, reordered or removed.\par\medskip
\noindent\textit{Editorial note.} Following the multi-claim convention for this source, source\_url is the source-level URL https://stacks.math.columbia.edu/ and the source label lemma-scheme-theoretic-support is carried in claim\_key and here rather than in a per-claim tag URL.\par\medskip
\noindent\textit{Editorial note.} The preamble is the chapter's shared Stacks Project preamble, retained as sources/stacks-morphisms/retained-preamble.tex: a line-for-line excerpt of the shared Stacks preamble staged with this snapshot. No macro was invented or redefined, and this card adds no preamble text of its own.\par\medskip
\noindent\textit{Editorial note.} Statement span 728-744 and proof span 746-769 come from the halfA pool card operations/pool-cards/stacks-morphisms-halfA.json and were re-read in the pinned chapter file before this card was written.\par\medskip
\noindent\textit{Editorial note.} The proof ends with the sentence 'We omit the proof.' after naming the algebra fact that the annihilator of a finite module commutes with localization at an element. That sentence is kept verbatim; what it leaves to the reader is that single elementary statement, which is displayed in the counted span, while the existence, uniqueness, minimality, finiteness and support assertions are all reduced inside the span.\par\medskip
\noindent\textit{Editorial note.} 5 cross-references occur inside the counted spans; they are listed in external\_references as navigation pointers only.\par\medskip
\noindent\textit{Allowed standard prerequisites (editorial).}\begin{itemize}
\item The same-chapter equivalence between quasi-coherent modules on a closed subscheme and quasi-coherent modules on the ambient scheme annihilated by the ideal sheaf, which gives the uniqueness in the form used here. Cited at proof lines 964-967 and again at 969-970. Converts the condition that the module is pushed forward from a closed subscheme into a statement about the annihilator, and gives uniqueness of the pushed forward module.
\item The same-chapter bijection between closed subschemes of an affine scheme and ideals of the coordinate ring, used to see that the annihilator rule defines a closed subscheme. Cited at proof lines 979-980. Reduces the existence of the smallest closed subscheme to showing that the annihilator rule gives a quasi-coherent sheaf of ideals.
\item Over an affine scheme a finite type quasi-coherent module is the tilde of a finite module, and the image of the support of the module is the closed subset defined by the annihilator. Cited at proof lines 971-978. Gives the finiteness of the pushed forward module and the identification of supports.
\end{itemize}\medskip
\ifdefined\hypertarget\hypertarget{sourcepointer-1}{}\fi
\noindent\textit{Source reference.} Original lemma 0560: focused-seeds/AGP/sources/stacks-snapshot-20261011/tags.txt, source lines 6693--6693. Same-chapter reference used in the retained author text at source line(s) 756. It is supplied only as a navigation pointer so that the retained text compiles; the counted item delivers the author's own statement and the author's complete proof, and no part of the referenced item is copied in or substituted for the counted argument.\par\medskip
\ifdefined\hypertarget\hypertarget{sourcepointer-2}{}\fi
\noindent\textit{Source reference.} Original lemma 00L2: focused-seeds/AGP/sources/stacks-snapshot-20261011/tags.txt, source lines 772--772. Same-chapter reference used in the retained author text at source line(s) 705, 759, 11631, 11640, 11679, 11688. It is supplied only as a navigation pointer so that the retained text compiles; the counted item delivers the author's own statement and the author's complete proof, and no part of the referenced item is copied in or substituted for the counted argument.\par\medskip
\ifdefined\hypertarget\hypertarget{sourcepointer-3}{}\fi
\noindent\textit{Source reference.} Original lemma 01QQ: focused-seeds/AGP/sources/stacks-snapshot-20261011/tags.txt, source lines 2275--2275. Same-chapter reference used in the retained author text at source line(s) 268, 511, 764, 815. It is supplied only as a navigation pointer so that the retained text compiles; the counted item delivers the author's own statement and the author's complete proof, and no part of the referenced item is copied in or substituted for the counted argument.\par\medskip
\ifdefined\hypertarget\hypertarget{sourcepointer-4}{}\fi
\noindent\textit{Source reference.} Original lemma 01QY: focused-seeds/AGP/sources/stacks-snapshot-20261011/tags.txt, source lines 2283--2283. Same-chapter reference used in the retained author text at source line(s) 578, 749, 782, 6149, 6758. It is supplied only as a navigation pointer so that the retained text compiles; the counted item delivers the author's own statement and the author's complete proof, and no part of the referenced item is copied in or substituted for the counted argument.\par\medskip
\ifdefined\hypertarget\hypertarget{sourcepointer-5}{}\fi
\noindent\textit{Source reference.} Original lemma 01PB: focused-seeds/AGP/sources/stacks-snapshot-20261011/tags.txt, source lines 2224--2224. Same-chapter reference used in the retained author text at source line(s) 703, 754, 4224, 5201, 5286, 6703, 10562, 11247. It is supplied only as a navigation pointer so that the retained text compiles; the counted item delivers the author's own statement and the author's complete proof, and no part of the referenced item is copied in or substituted for the counted argument.\par\medskip
\begin{lemma}
\label{lemma-scheme-theoretic-support}
Let $\mathcal{F}$ be a finite type quasi-coherent module
on a scheme $X$. There exists a smallest closed subscheme
$i : Z \to X$ such that there exists a quasi-coherent
$\mathcal{O}_Z$-module $\mathcal{G}$ with
$i_*\mathcal{G} \cong \mathcal{F}$. Moreover:
\begin{enumerate}
\item If $\Spec(A) \subset X$ is any affine open, and
$\mathcal{F}|_{\Spec(A)} = \widetilde{M}$ then
$Z \cap \Spec(A) = \Spec(A/I)$ where $I = \text{Ann}_A(M)$.
\item The quasi-coherent sheaf $\mathcal{G}$ is unique up to unique
isomorphism.
\item The quasi-coherent sheaf $\mathcal{G}$ is of finite type.
\item The support of $\mathcal{G}$ and of $\mathcal{F}$ is $Z$.
\end{enumerate}
\end{lemma}

\begin{proof}
Suppose that $i' : Z' \to X$ is a closed subscheme which satisfies the
description on open affines from the lemma. Then by
Lemma \ref{lemma-i-star-equivalence}
we see that $\mathcal{F} \cong i'_*\mathcal{G}'$ for some unique
quasi-coherent sheaf $\mathcal{G}'$ on $Z'$. Furthermore, it is clear
that $Z'$ is the smallest closed subscheme with this property (by the
same lemma). Finally, using
Properties, Lemma \ref{properties-lemma-finite-type-module}
and
Algebra, Lemma \ref{algebra-lemma-finite-over-subring}
it follows that $\mathcal{G}'$ is of finite type. We have
$\text{Supp}(\mathcal{G}') = Z$ by
Algebra, Lemma \ref{algebra-lemma-support-closed}.
Hence, in order to prove the lemma it suffices to show that
the characterization in (1) actually does define a closed subscheme.
And, in order to do this it suffices to prove that the given rule
produces a quasi-coherent sheaf of ideals, see
Lemma \ref{lemma-closed-immersion-bijection-ideals}.
This comes down to the following algebra fact: If $A$ is a ring, $f \in A$,
and $M$ is a finite $A$-module, then
$\text{Ann}_A(M)_f = \text{Ann}_{A_f}(M_f)$.
We omit the proof.
\end{proof}

\end{document}
